\documentstyle[%


righttag%


,11pt%


,amssymb%


,russian%               remove in Eng.


,izvru%                 Set izven% in Eng.


]{amsart}





% 





\newcommand{\vfi}{\varphi}


\newcommand{\al}{\alpha}


\newcommand{\ls}{\leqslant}


\newcommand{\rs}{\geqslant}


\newcommand{\R}{\Bbb R}


\newcommand{\de}{\doteq}


\newcommand{\fy}{\infty}


\newcommand{\Pt}{\partial}


\newcommand{\ve}{\varepsilon}


\newcommand{\lr}{\longrightarrow}


\newcommand{\intr}{\operatorname{int}}


\newcommand{\rank}{\operatorname{rank}}


\newcommand{\tts}[1]{}





%\hoffset=-15mm%                  For Eng.


 


\setcounter{page}{68}


\god{1999}


\nomer{2~(441)} %                        Remove in Eng.


%\nomer{Vol.\,43, No.\,2}%               Set in Eng.


%\str{\pageref{begin}--\pageref{end}}%  Set in Eng.


\udk{517.934}





\author{\it Ю.В.\,Мастерков}


% NB:   ^^ remove in Eng.


\title{К вопросу о локальной управляемости\\


в критическом случае}


\thanks{Работа поддержана Российским фондом


фундаментальных исследований (грант 97--01--00413) и конкурсным центром


Удмуртского госуниверситета (грант 97--04).}





\date{}


%\pagestyle{myheadings}%         Set in Eng.


\pagestyle{plain} %              Remove in Eng.


\begin{document}


%\thispagestyle{plain}%          Set in Eng.


\maketitle


\label{begin}





\section*{Введение}





Здесь изучаются условия управляемости в нуль системы


\begin{equation}


\Dot x = f_0(x) + uf_1(x),\quad (x,u)\in \R^n\times[-1,1]


\end{equation}


в критическом случае (т.\,е.~в случае, когда система линейного приближения


не является вполне управляемой). В работе введено понятие устойчивой


локальной управляемости и показано, что если система линейного приближения


вполне управляема, то система (1) устойчиво управляема; если же свойство


вполне управляемости системы линейного приближения отсутствует, то реализуется


один из трех случаев: 1) система (1) устойчиво локально управляема;


2) система (1) управляема; 3) система (1) неуправляема. Выделен класс


систем вида (1), для которых имеет место устойчивая локальная управляемость.


Показано, что $N$-управляемость [1] системы


\begin{equation}


\Dot x = f(x,u),\quad x\in\R^n,\quad u\in U\subset\R^m


\end{equation}


влечет


устойчивую управляемость. Таким образом, свойство $N$-управляемости


является наиболее сильным из известных к настоящему времени свойств


управляемости, а устойчивая управляемость занимает промежуточное


положение между $N$-управляемостью и локальной управляемостью.





Задачам исследования условий управляемости систем вида (2) посвящен ряд


работ [1]--[10], причем в [2], [7], [10] получены достаточные


условия управляемости в критическом случае. Полный ответ на вопрос об


устойчивой управляемости получен лишь для систем второго порядка. Для


произвольного $n>2$ полный ответ о локальной управляемости остается


открытым. В некоторых работах проверка управляемости в критическом


случае предполагает наличие некоторых объектов, существование и


построение которых весьма проблематично [10]. Здесь возникает ситуация,


аналогичная построению функций Ляпунова во втором методе Ляпунова, что


не умаляет достоинств таких утверждений.





\section{Обозначения и определения}





В работе используются следующие


обозначения: $\R^n$ --- евклидово пространство размерности $n$ с нормой


$|x|=\sqrt{x^*x}$ ($*$ --- операция транспонирования); $\R\de\R^1$;


$\R^+\de [0,\fy)$; $O^n_{\ve}(x_0)\de\{x\in\R^n : |x-x_0|<\ve\}$;


$O^n_\ve\doteq O^n_{\ve}(0)$; $\intr G$ --- внутренность, $\ov G$


--- замыкание, $\Pt G$ --- граница множества $G\subset\R^n$;


$C^n(\R^k,\R^p)$ --- пространство $n$ раз непрерывно-дифференцируемых


функций из $\R^k$ в $\R^p$.





Рассмотрим систему (2) в предположении, что $U$ --- непустой выпуклый


компакт в $\R^n$, $0\in\intr U$, $f\in C^1(\R^n\times\R^m,\R^n)$,


$f(0,0)=0$ и матрица $\left(\frac{\Pt f(x,u)}{\Pt u}\right)_{(0,0)}$


ненулевая.





Всевозможные измеримые вектор-функции $u : \R\lr U$ будем называть


{\it допустимыми управлениями}.





Точка $x_0$ называется $\tau$-управляемой, если существует такое


допустимое управление $u:[0,\tau]\lr U$, что разрешима краевая задача


$\Dot x = f(x,u(t))$, $x(0)=x_0$, $x(\tau)=0$.





Множество всех $\tau$-управляемых точек называется {\it множеством


управляемости} системы (2) {\it за время} $\tau$ и обозначается


$D_{\tau}$.  Известно (напр., [11], гл.\,1, \S\,2, с.\,78), что при


всех $\tau>0$ множество $D_{\tau}$ --- компакт в $\R^n$.  Множество


$D_{\fy}=\underset{\tau>0}{\cup}{}D_{\tau}$ называется {\it множеством


управляемости} системы (2).





Напомним, что система (2) называется локально управляемой (напр.,


[11], гл.\,1, \S\,2, с.\,77; [12], гл.\,1, \S\,7, с.\,60), если $0\in\intr


D_{\tau}$ при некотором $\tau>0$.





\begin{defn}


Сиcтема (2) называется устойчиво управляемой, если


для любого $\ve>0$ существует такое $\delta>0$, что для каждого $x_0\in


O_{\delta}^n$ найдутся время $\tau=\tau_{x_0}>0$ и такое допустимое


управление $u_{x_0}:[0,\tau]\lr U$, что решение


$x(t)$, $t\in[0,\tau]$, системы (2) при данном управлении $u=u_{x_0}(t)$


и начальном условии $x(0)=x_0$ удовлетворяет условиям: $x(\tau)=0$,


$|x(t)|<\ve$, $t\in[0,\tau]$.


\end{defn}





В [1] введено понятие $N$-{\it управляемости}, которое


означает, что $0\in\intr D_{\tau}$ при всех $\tau>0$.





Пусть $\frak L$, $\frak M$, $\frak N$ соответственно множества локально


управляемых, устойчиво управляемых и $N$-управляемых систем вида (2).





\begin{lem}


$\frak N\subset\frak M\subset\frak L$.


\end{lem}





\begin{pf}


Докажем, что $\frak N\subset\frak M$. Действительно, если


система (2) $N$-управляема, то для любого $\ve>0$ найдется такое


$\tau>0$, что $D_{\tau}\subset O_{\ve}^n$. Пусть


$\delta=\min\limits_{x\in\Pt D_{\tau}} |x|$. Тогда существует такое решение


$x(t)$ системы (2), что $x(0)=x_0$, $x(\tau)=0$,


$x(t)\subset D_{\tau}\subset O_{\ve}^{\,n}$ при всех $t\in[0,\tau]$,


т.\,е.~система (2) устойчиво управляема.





Включение $\frak M\subset\frak L$ очевидно.


\end{pf}





Приводимые ниже примеры показывают, что множества


$\frak N$, $\frak M$ и $\frak L$ не совпадают.





\begin{exam}


Рассмотрим систему


\begin{equation}


\Dot x_1=u,\quad


\Dot x_2=


\frac{1}{2}(5x_1^4-8x_1^3+3x_1^2)+\frac{1}{2}(5x_1^4-8x_1^3+3x_1^2)u+u,


\end{equation}


где $u\in[-1,1]$. Здесь (см. (1))


$f_0(x)=(0,\frac{1}{2}(5x_1^4-8x_1^3+3x_1^2))^{*}$;


$f_1(x)=(1,\frac{1}{2}(5x_1^4-8x_1^3+3x_1^2))^{*}$;


$b=f_1(0)=(1,1)^{*}$.


Очевидно, $A=\left(\frac{\Pt f_0(x)}{\Pt x}\right)_{x=0}$ --- нулевая


матрица и, следовательно, $\rank (b,Ab)=1$, т.\,е.~к системе (3) не


применима теорема о локальной управляемости по первому приближению. Тем


не менее, система (3) является локально управляемой. Действительно,


кривые $x_2=x_1^5-2x_1^4+x_1^3+x_1+C_1$ являются траекториями решений


системы (3) при $u=1$. Движение по ним осуществляется в сторону


возрастания $x_1$. Соответственно прямые $x_2=x_1+C_2$ суть


траектории допустимых решений системы (3) при $u=-1$ и движение по ним


осуществляется в сторону убывания $x_1$. Пусть


$\gamma_{+}\de\{x_2=x_1^5-2x_1^4+x_1^3+x_1,\,x_1\ls 0\}$,


$\gamma_{-}\de\{x_2=x_1,\,x_1\rs 0\}$ (см. рис.~1). Легко видеть, что из


любой точки $\beta\in\R^2$, находящейся ниже кривой


$\gamma=\gamma_{+}\cup\gamma_{-}$, двигаясь по прямой $x_2=x_1+C(\beta)$,


изображающая точка системы (3) за конечное время достигает траектории


$\gamma_{+}$ и далее, двигаясь по ней, за конечное же время достигает


нуля. Как осуществляется переход в нуль из точек, находящихся выше


кривой $\gamma$, поясним по рисунку~1, на котором движение в


заштрихованной области показано справа в более крупном масштабе. Из


произвольной точки $\al\in\R^2$, находящейся выше кривой $\gamma$, по


траектории $\{x_2=x_1^5-2x_1^4+x_1^3+x_1+C(\al)\}$ за конечное время


можно достичь полосы $0.6\ls x_1\ls 1$, в которой (см. рис.~1


справа) осуществляется спуск до траектории $\gamma_{-}$ и далее, двигаясь


по $\gamma_{-}$, можно достичь начала координат за конечное время.


Таким образом, система (3) локально управляема (отметим, что для


оптимального по быстродействию перехода из изображенной на


рис.~1 точки $\al$ в нуль потребовалось 27 переключений управления).


\begin{center}


\unitlength=1mm


\begin{picture}(170,95) %% width, heigth, added 5 mm for signature


\put(19,95){\special {em:graph m-12.bmp}} %% left X, upper Y, -, /2, +toX


\put(78,0){\mbox Рис.\,1}


\end{picture}


\end{center}


Покажем, что система (3) не является устойчиво управляемой. Действительно,


при $\ve=0.6$ никакую точку $x_0\in O^n_{\ve}$, находящуюся выше кривой


$\gamma$, нельзя перевести в нуль за конечное время, не покидая при этом


окрестности $O^n_{\ve}$. Дело в том, что в $O^n_{\ve}$ поле скоростей


системы (3) направлено вверх относительно любой опорной прямой к


кривой $\gamma$. Следовательно, никакая допустимая траектория в


$O^n_{\ve}$ не сможет подойти ``сверху'' к траектории $\gamma$, а


значит, и попасть в начало координат, не покидая при этом окрестности


$O^n_{\ve}$, т.\,е.~система (3) не является устойчиво управляемой.


\end{exam}





Покажем, что множества $\frak N$ и $\frak M$ не совпадают.





\begin{exam}





Рассмотрим систему





% \begin{center}


% \begin{picture}(470,185)


%   \begin{picture}(200,185)(12,-85) % табулятор


%     \begin{tabular}{p{8cm}}





\begin{equation}


\Dot x_1=x_2+u(1-x_2),\quad


\Dot x_2=-x_1+ux_1,


\end{equation}


где $u\in[0,1]$. Нетрудно заметить, что окружности $x_1^2+x_2^2=C_1^2$ и


прямые $x_2=C_2$ суть траектории решений системы (4) при $u=0$ и $u=1$


соответственно, причем движение по окружностям осуществляется по


часовой стрелке, а по


% \end{tabular}


%  \end{picture}


%  \end{picture}


%\end{center}


прямым --- в сторону возрастания $x_1$ (см. рис.~2).





Покажем, что система (4) устойчиво управляема. Действительно, пусть


$\ve$ --- произвольное положительное число, а $x_0$ --- произвольная


точка, удовлетворяющая условию $|x_0|<\ve$.


Двигаясь по окружности $x_1^2+x_2^2=|x_0|^2$, за


время $t_1<2\pi$ можно достичь отрицательной полуоси $Ox_1$, и далее,


двигаясь по траектории $\gamma_{+}\de\{x_2=0, x_1\ls 0\}$, за время


$t_2=|x_0|$ достигаем начала координат, оставаясь все время в


окрестности $O^n_{\ve}$. Однако при всех $\tau\ls\pi/4$


множество $D_{\tau}$ содержит нуль на своей границе, т.\,е.~система (4)


не является $N$-управляемой.  Отметим, что приблизиться к нулю можно


лишь со стороны третьего квадранта. Чтобы достичь нуля из точек


положительной полуоси $Ox_1$, необходимо сначала достичь третьего


квадранта, затратив на это время $\tau=\pi/4$, следовательно, при


всех $\tau\ls\pi/4$ нуль не принадлежит внутренности $D_{\tau}$,


т.\,е.~система (4) не является $N$-управляемой.


\begin{center}


\unitlength=1mm


\begin{picture}(170,60) %% width, heigth, added 5 mm for signature


\put(51,60){\special {em:graph m-3.bmp}} %% left X, upper Y, -, /2, +toX


\put(78,0){\mbox Рис.\,2}


\end{picture}


\end{center}





\end{exam}





Вообще говоря, система (4) не является системой вида (2), т.\,к.~нуль


находится на границе множества управляемости $U=[0,1]$. Корректный


пример, показывающий, что при отсутствии $N$-управляемости может иметь


место устойчивая управляемость, легко построить на основе примера 2.





Рассмотрим систему


\begin{equation}


\Dot x_1=x_2+u_1^2(1-x_2),\quad


\Dot x_2=-x_1+u_1^2x_1,\quad


\Dot x_3=u_2,


\end{equation}


где $u_{i}\in[-1,1]$, $i=1,2$. Так как траектории системы (4) суть проекции


траекторий системы (5), то система (5) не является $N$-управляемой.


Устойчивая управляемость следует из того, что из любой точки


$x_0=(x_{01},x_{02},x_{03})^{*}$ при $u_2=0$ можно достичь точки


$x_{03}$ на оси $ox_3$, а затем при $u_2={\mathrm{sign}}(-x_{03})$,


$u_1=0$ достичь нуля, не покидая при этом окрестности $O^n_{|x_0|}$


(решение здесь понимается в смысле А.Ф.\,Филиппова).





\section{Постановка задачи. Основные утверждения}





Рассмотрим систему


\begin{equation}


\Dot x = f_0(x) + uf_1(x),\quad (x,u)\in \R^n\times[-1,1].


\end{equation}


Предполагается, что $f_{i}\in C^p(\R^n,\R^n)$, $i=0,1$, $f_0(0)=0$,


$f_1(0)=b\ne 0$. Наряду с системой (6) рассмотрим соответствующую


ей систему линейного приближения


\begin{equation}


\Dot x = Ax + bu,


\end{equation}


где $A=\left(\frac{\Pt f_0(x)}{\Pt x}\right)_{x=0}$. Известны следующие


утверждения (напр., [11], гл.\,1, \S\,3, с.\,91; [12], гл.\,1, \S\,7,


с.\,61).





\begin{lem}


Предположим, что $\rank(b,Ab,A^2b,\dotsc,A^{n-1}b)=k\ls n$,


тогда векторы $b$, $Ab$, $A^2b,\dotsc,A^{k-1}b$ линейно независимы.


\end{lem}





\begin{lem}


Пусть


$\rank(b,Ab,A^2b,\dotsc,A^{n-1}b)=k\ls n$ и $L^k$ --- линейное


подпространство в $\R^n$, порожденное векторами


$b$, $Ab$, $A^2b,\dotsc,A^{k-1}b$. Тогда множество управляемости


$D_{\fy}$ системы $(7)$ содержится в $L^k$.


\end{lem}





\begin{note}


В условиях леммы 3 система (7) является $N$-управляемой в


$L^k$ (напр., [11], гл.\,2, \S\,3, с.\,108).


\end{note}





\begin{lem}


Пусть $\rank(b,Ab,A^2b,\dotsc,A^{n-1}b)=n-1$,


тогда для любого $\tau>0$ существует


гладкое многообразие


$M^{n-1}$, которое содержится в множестве управляемости $D_{\tau}$


системы $(6)$.


\end{lem}





\begin{pf}


В силу замечания к лемме 3


система (7) является $N$-управляемой на многообразии $L^{n-1}$,


порожденном векторами $b$, $Ab$, $A^2b,\dotsc,A^{n-2}b$.


Следовательно, для любого $\tau>0$ существуют


допустимые управления $u_i(t)$, $t\in[0,\tau]$, $i=\ov{1,n-1}$,


которые переводят систему с обращенным временем


\begin{equation}


\Dot x = -Ax - bu


\end{equation}


из начала координат в линейно независимые


точки $x_i(\tau)=x(\tau,u_i(\cdot))$, $i=\ov{1,n-1}$.


Без ограничения общности будем полагать, что:


а)~$u_i(t)\in C^{\fy}([0,\tau])$, $i=\ov{1,n-1}$;


б)~управления $u_i(t)$ выбраны так, что соответствующие им


решения $x_i(t)=x(t,u_i(\cdot))$ системы (8)


существуют на отрезке $[0,\tau]$ и удовлетворяют условию


$x_i(1)=\ve A^{i-1}b$, $i=\ov{1,n-1}$, где $\ve$ --- некоторое


положительное число. Обозначим


$u(t,\xi)=\sum\limits_{i=1}^{n-1}\xi_i u_i(t)$,


где $|\xi|\ls 1$.





Пусть $x(t,\xi)$ --- решение системы с обращенным временем


$\Dot x = -f_0(x) - u(t,\xi)f_1(x)$, удовлетворяющее начальному


условию  $x(0,\xi) = 0$.


Можно полагать, что при всех


$\xi\in O^{n-1}_1$ данное решение существует на отрезке $[0,\tau]$.


Обозначим


$Z(t)=\left(\frac{\Pt x(t,\xi)}{\Pt\xi}\right)_{\xi=0}$.


Легко видеть, что $n\times(n-1)$-матрица $Z(t)$


удовлетворяет уравнению


$$


\Dot Z(t) = -AZ(t) - bu^{*}(t),


$$


где $bu^{*}(t)$ --- матричное произведение


вектор-столбца $b$ на вектор-строку


$u^{*}(t)=(u_1(t),u_2(t),\linebreak \dotsc,u_{n-1}(t))^{*}$.


Отсюда следует, что столбцы


$z_i(t)$, $i=\ov{1,n-1}$, матрицы $Z(t)$


удовлетворяют уравнениям


$$


\Dot z_i(t)=-Az_i(t)-bu_i(t),\quad z_i(0)=0,\quad i=\ov{1,n-1}.


$$


Но тогда $z_1(\tau),z_2(\tau),\dotsc,z_{n-1}(\tau)$ ---


линейно независимые векторы.


Следовательно, отображение $\xi\lr x(\tau,\xi)$


является непрерывно-дифференцируемым отображением


окрестности $O^{n-1}_1$ в $\R^n$.


Заметим теперь, что поскольку


$\rank Z(\tau)=\rank\left(\frac{\Pt x(\tau,\xi)}{\Pt\xi}\right)_{\xi=0}=n-1$,


то существует такая окрестность $O^{n-1}_{\delta}$, что


$\rank\left(\frac{\Pt x(\tau,\xi)}{\Pt\xi}\right)=n-1$ для всех $\xi\in


O^{n-1}_{\delta}$.  А это означает, что множество


$M^{n-1}_{\tau}\de\{x=x(\tau,\xi),\,\xi\in O^{n-1}_{\delta}\}$ будет


гладким многообразием, причем $M^{n-1}_{\tau}\subset D_{\tau}$.


\end{pf}





\begin{cor}


Пусть $\rank(b,Ab,A^2b,\dotsc,A^{n-1}b)=n-1$.


Тогда для любого $\tau>0$ существуют окрестность


$O^n_{\ve}$ и такая функция $\vfi_{\tau}\in C^1(\R^n,\R)$,


что $(\Pt\vfi_{\tau}(x)/\Pt x)\ne 0$ для всех $x\in O^n_{\ve}$


и $M^{n-1}\de\{x\in O^n_{\ve} : \vfi_{\tau}(x)=0\}\subset D_{\tau}$.


\end{cor}





\begin{pf}


Из доказательства леммы 4 видно, что для любого


$\tau>0$ существует окрестность $O^{n-1}_{\delta}$, для которой многообразие


$M^{n-1}_{\tau}\de\{x=x(\tau,\xi),\,\xi\in O^{n-1}_{\delta}\}$ содержится


в $D_{\tau}$, причем $\rank\left(\frac{\Pt x(\tau,\xi)}{\Pt\xi}\right)=n-1$


для всех $\xi\in O^{n-1}_{\delta}$. Известно (напр., [13], гл.\,7,


\S\,7, утв.\,1, с.\,505), что в этом случае существуют такая окрестность


$O^n_{\ve}$, индекс $k$ и функция


$\vfi^k_{\tau}(x_1,\dotsc,x_{k-1},x_{k+1},\dotsc,x_n)$, что


$M^{n-1}_{\tau}\de\{x\in O^n_{\ve}:


x_k=\vfi^k_{\tau}(x_1,\dotsc,x_{k-1},x_{k+1},\dotsc,x_n)\}$.


Это равносильно тому, что $M^{n-1}\de\{x\in


O^n_{\ve}:\vfi_{\tau}(x)=0\}$, где


$\vfi_{\tau}(x)=x_k-\vfi^k_{\tau}(x_1,\dotsc,x_{k-1},x_{k+1},\dotsc,x_n)$.


Заметим теперь, что $(\Pt\vfi_{\tau}(x)/\Pt x)\ne 0$


для всех $x\in O^n_{\ve}$.


\end{pf}





\begin{note}


На самом деле в условиях леммы 4


гладкость многообразия $M^{n-1}_{\tau}$ будет степени $p$. Это следует


из теоремы о степени гладкости зависимости решений от параметра


(напр., [14], гл.\,3, \S\,21, с.\,88).


\end{note}





Пусть


$\rank(b,Ab,A^2b,\dotsc,A^{n-1}b)=n-1$.


Рассмотрим построенное в доказательстве леммы 4 многообразие


$M^{n-1}_{\tau}\de\{x=x(\tau,\xi),\,\xi\in O^{n-1}_{\delta}\}\subset


D_{\tau}$. Векторы $p_i(x_0)\de\left(\frac{\Pt


x(\tau,\xi)}{\Pt\xi_i}\right)_{\xi=\xi_0}$,


$i=\ov {1,n-1}$, очевидно, образуют


базис касательного пространства к многообразию $M^{n-1}_{\tau}$ в точке


$x_0=x(\tau,\xi_0)$.  Введем функцию


$s(x)=\det (f_0(x),p_1(x),\dotsc,p_{n-1}(x))$.





\begin{lem}


Если в любой окрестности $O_{\ve}^n$ найдутся


такие точки $x_{+}$ и $x_{-}$, принадлежащие $M^{n-1}_{\tau}$,


что $s(x_{+})s(x_{-})<0$, то система


$(6)$ устойчиво управляема.


\end{lem}





\begin{pf}


Зафиксируем $\ve>0$. Легко видеть, что


существует такое $\ve_1>0$, что любая точка


$x_0\in O_{\ve_1}^n\cap M^{n-1}_{\tau}$ под действием допустимого


управления $u_{x_0}=u(t,\xi_{x_0})$ может быть переведена в нуль


за конечное время, не покидая при этом окрестности $O_{\ve}^n$.


Это следует из теоремы о непрерывной зависимости от параметра


(напр., [14], гл.\,3, \S\,20, стр.\,86).





Далее поскольку


$M^{n-1}_{\tau}\de\{x=x(\tau,\xi),\,\xi\in O^{n-1}_{\delta}\}$,


то, не нарушая общности, можно полагать, что многообразие


$M^{n-1}_{\tau}$ делит окрестность


$O_{\theta}^n=O_{\ve_1}^n\cap O_{\delta}^n$


на две части $O_{\theta}^{+}$


и $O_{\theta}^{-}$, или ``верхнюю'' и ``нижнюю''


соответственно. Положим для определенности, что $s(x_{+})>0$.  Это


означает, что вектор $f_0(x_{+})$ направлен трансверсально к


многообразию $M^{n-1}_{\tau}$ в точке $x_{+}$.


Будем считать, что он направлен ``вверх'' относительно


$M^{n-1}_{\tau}$. Тогда $s(x_{-})<0$ и, следовательно, вектор


$f_0(x_{-})$ направлен ``вниз'' относительно


$M^{n-1}_{\tau}$.  Следовательно, существует такое $t_0>0$, что


$x(t,x_{+})\in O_{\theta}^{+}$, а $x(t,x_{-})\in O_{\theta}^{-}$ для всех


$t\in(0,t_0)$, где $x(t,x_0)$ --- решение системы $\Dot x = -f_0(x)$,


удовлетворяющее начальному условию $x(0,x_0)=x_0$.





Легко видеть, что любая точка $x(s)_{+}=x(s,x_{+})$,


$x(s)_{-}=x(s,x_{-})$, $s\in[0,t_0]$,


под действием допустимого


управления может быть переведена в нуль


за конечное время, не покидая при этом окрестности $O_{\ve}^n$.





Заметим, что найдется такая окрестность $O_{\theta_1}^n$,


что множество решений $x=x(t,\xi,x(s)_{+})$


системы с обращенным временем $\Dot x = -f_0(x) - u(t,\xi)f_1(x)$,


удовлетворяющих начальному условию


$x(0,\xi)=x(s)_{+}$, целиком покрывает множество


$O_{\theta_1}^n\cap O_{\theta}^{+}$.


А это означает, что любая точка $x_0\in O_{\theta_1}^n\cap O_{\theta}^{+}$


может быть переведена в точку $x(s)_{+}$ за конечное время,


не покидая при этом окрестности $O_{\ve}^n$.


Аналогично показывается, что существует такая окрестность


$O_{\theta_2}^n$, что любая точка $x_0\in O_{\theta_2}^n\cap O_{\theta}^{-}$


может быть переведена в точку $x(s)_{-}$ за конечное время,


не покидая при этом окрестности $O_{\ve}^n$.


Пусть $O_{\theta_0}^n=O_{\theta_1}^n\cap O_{\theta_2}^n$. Тогда получаем, что


из любой точки $x_0\in O_{\theta_0}^n$ за конечное время можно


достичь точки $x(s)_{+}$ или $x(s)_{-}$, а затем, двигаясь по


допустимым траекториям системы (6), за конечное же время достичь


нуля, не покидая при этом окрестности $O_{\ve}^n$. Значит, система


(6) устойчиво управляема.


\end{pf}





\begin{thm}


Предположим, что


$\rank(b,Ab,A^2b,\dotsc,A^{n-1}b)=n-1$. Пусть далее


$\gamma_{\tau}(\xi_1)=x(\tau,\xi_1)$, $\xi_1\in[-1,1]$, где $x(t,\xi_1)$


--- решение системы


\begin{equation}


\Dot x=f_0(x)+\xi_1f_1(x),\quad x(0,\xi)=0.


\end{equation}


Для того чтобы система $(6)$ была устойчиво управляемой,


достаточно, чтобы функция $s(\xi_1)=s(\gamma_{\tau}(\xi_1))$ меняла в


нуле знак.


\end{thm}





{\bf Доказательство}. Так как $\gamma_{\tau}(\xi_1)\in M^{n-1}_{\tau}$,


то для доказательства достаточно сослаться на лемму 5.





\begin{thm}


Пусть существует функция $\vfi\in C^1(\R^n,\R)$,


удовлетворяющая условию$:$ всякому $\tau>0$ отвечает такое $\ve>0$,


что $(\Pt\vfi(x)/\Pt x)\ne 0$ для всех $x\in


O^n_{\ve}$ и $M\de\{x\in O^n_{\ve}:\vfi(x)=0\}\subset D_{\tau}$.





Если существуют $\tau_0\in(0,\tau)$ и такие решения


$x_1(t)=x(t,u_1(\cdot))$, $x_2(t)=x(t,u_2(\cdot))$ системы $(6)$,


что $\vfi(x_1(t))\vfi(x_2(t))<0$ для всех $t\in(0,\tau_0]$, то


система $(6)$ $N$-управляема.


\end{thm}





\begin{pf}


Обозначим через $x(t,u_i(\cdot),x_0)$ решения системы


(6), соответствующие управлению $u_i(t)$, $t\in[0,\tau_0]$, и начальному


условию $x(t,u_i(\cdot),x_0)=x_0$, $i=1,2$. Отметим, что


множество $M$ является $(n-1)$-мерным гладким многообразием,


разделяющим окрестность $O^n_{\ve}$ на две части:


$O_{\ve}^{+}\de\{x\in O^n_{\ve}:\vfi(x)>0\}$ и


$O_{\ve}^{-}\de\{x\in O^n_{\ve}:\vfi(x)<0\}$. Тогда неравенство


$\vfi(x_1(t))\vfi(x_2(t))<0$ означает, что одно из решений $x_i(t)$,


$i=1,2$, лежит в $O_{\ve}^{+}$, а другое соответственно в $O_{\ve}^{-}$


для всех $t\in(0,\tau_0]$.


Пусть для определенности $x_1(t)\in O_{\ve}^{+}$, а $x_2(t)\in O_{\ve}^{-}$.


Тогда для любого $t_0\in(0,\tau_0)$ существует такая окрестность


$O^n_{\delta}$, что $x(t_0,u_1(\cdot),x_0)\in O_{\ve}^{+}$,


а $x(t_0,u_2(\cdot),x_0)\in O_{\ve}^{-}$ для всех $x_0\in O^n_{\delta}$.


Таким образом, $\vfi(x(0,u_1(\cdot),x_0)<0$,


$\vfi(x(t_0,u_1(\cdot),x_0)>0$ для каждого


$x_0\in O^n_{\delta}\cap O^{-}_{\ve}$. И, поскольку функция


$\vfi(x(t,u_1(\cdot),x_0)$ непрерывна, существует такое время


$t_{x_0}\in(0,t_0)$, что $\vfi(x(t_{x_0},u_1(\cdot),x_0)=0$,


т.\,е.~$x(t_{x_0},u_1(\cdot),x_0)\in M\subset D_{\tau}$. Аналогично,


для каждой точки $x_0\in O^n_{\delta}\cap O^{+}_{\ve}$ существует такое


$t_{x_0}\in(0,t_0)$, что $x(t_{x_0},u_2(\cdot),x_0)\in M\subset D_{\tau}$.


Следовательно, из любой точки $x_0\in O^n_{\delta}$ за время $t\ls t_0$


можно попасть на многообразие $M\subset D_{\tau}$ и затем за время $\tau$


в начало координат, не покидая при этом множества $D_{\tau+t_0}$.


И поскольку $\tau$ и $t_0$ выбраны произвольно, а $D_{\tau+t_0}$


--- компакт, то система $N$-управляема.


\end{pf}





Повторением основных моментов доказательства леммы 5 и теоремы 2


обосновывается





\begin{thm}


Пусть функция $\vfi\in C^1(\R^n,\R)$


удовлетворяет условиям: {\em а)}~существует такое $\ve>0$, что


$(\Pt\vfi(x)/\Pt x)\ne 0$ для всех $x\in O^n_{\ve};$


{\em б)}~существует такое $\tau>0$, что


$M\de\{x\in O^n_{\ve}:\vfi(x)=0\}\subset D_{\tau}$.





Если существуют такие моменты времени $t_i\in(0,\tau)$ и решения системы


$(6)$ $x_i(t)=x(t,u_i(\cdot))$, $t\in[0,\tau_0]$, $i=1,2$,


что $\vfi(x_1(t_1))\vfi(x(t_2))<0$, то


система $(6)$ локально управляема.


\end{thm}








\begin{thebibliography}{99}





\bibitem{1}


Петров Н.Н. {\em Локальная управляемость автономных


систем} // Дифференц. уравнения. -- 1968. -- Т.\,4. -- \No\,7. --


С.\,1218--1232.


\bibitem{2}


Петров Н.Н. {\em Об управляемости автономных


систем} // Дифференц. уравнения. -- 1968. -- Т.\,4. -- \No\,4. --


С.\,606--617.


\bibitem{3}


Петров Н.Н. {\em Решение одной задачи теории управляемости} //


Дифференц. уравнения. -- 1969. -- Т.\,5. -- \No\,5. -- С.\,962--963.


\bibitem{4}


Емельянов С.В., Коровин С.К.,


Мамедов И.Г., Никитин С.В. {\em


Критерии управляемости нелинейных систем при фазовых ограничениях} //


ДАН СССР. -- 1986. -- Т.\,290. -- \No\,1. -- С.\,18--22.


\bibitem{5}


Землякова Л.С. {\em Управляемость нелинейных систем


дифференциальных уравнений} // Дифференц. уравнения:


качеств. теория. -- Рязань, 1995. -- С.\,64--71.


\bibitem{6}


Копейкина Т.Б. {\em К необходимым условиям управляемости


нелинейных систем в критическом случае} // Препринт.


Ин-т матем. АН БССР. -- Минск, 1985. -- \No\,27. -- 44 с.


\bibitem{7}


Копейкина Т.Б. {\em О локальной управляемости


нелинейных систем в критическом случае} //


Весцi АН БССР. Сер. фiз.-матэм.


наук. -- 1987. -- \No\,1. -- С.\,8--15.


\bibitem{8}


Митрохин Ю.С., Степанов А.Н.


{\em Критические случаи управляемости систем нелинейных дифференциальных


уравнений оптимального регулирования} // Дифференц. уравнения:


качеств. теория. --  Рязань, 1985. -- С.\,61--70.


\bibitem{9}


Мастерков Ю.В. {\em Об устойчивой локальной


нуль-управляемости систем с квадратичной нелинейностью


на плоскости} // Изв. отд. матем. и информ. -- Ижевск, 1993. --


\No\,2.  -- С.\,3--24.


\bibitem{10}


Никольский М.С. {\em Об условиях второго порядка в задаче о


нуль управляемости} //


Дифференц. уравнения. -- 1968. -- Т.\,4 -- \No\,7. -- С.\,137.


\bibitem{11}


Ли Э.М., Маркус Л. {\em Основы теории оптимального


управления.} -- М.: Наука, 1972. -- 576 с.


\bibitem{12}


Габасов Р., Кириллова Ф.М.


{\em Качественная теория оптимальных процессов.} -- М.: Наука, 1971. -- 507 с.


\bibitem{13}


Зорич В.А. {\em Математический анализ}. Часть 1. Учеб. пособие. -- М.: Наука,


1981.  -- 543 с.


\bibitem{14}


Петровский И.Г. {\em Лекции по теории обыкновенных


дифференциальных уравнений.} -- М.: Изд-во МГУ, 1984. -- 295 с.





\end{thebibliography}





\vskip 2cm





\post{\it Удмуртский государственный}{\it Поступила}


\post{\it университет}{03.03.1998}





\label{end}


\end{document}





